\ifdefined\gkver\else\def\gkver{student}\fi
\documentclass[\gkver]{gaokaozhenti}
\begin{document}

\chapter{1990年全国}

\begin{enumerate}
\item
%% number: 1
%% paperName: 1990年全国高考第 1 题 2 分
%% typeId: 1
%% type: 单选题
%% chapter: 其他
%% point: 物理学史
%% method: 无
%% score: 2
%% degree: 950
%% duplicateId: 0
%% body:
第一个发现电磁感应现象的科学家是\xzanswer{C}
\fourchoices[answer=C]
{奥斯特}
{库仑}
{法拉第}
{安培}

%% answer: C
%% memo:
\memoanswer{
第一个发现电磁感应现象的科学家是法拉第，故 C 正确；奥斯特是第一个发现电流的磁效应的物理学家。
}

\item
%% number: 2
%% paperName: 1990年全国高考第 2 题 2 分
%% typeId: 1
%% type: 单选题
%% chapter: 机械振动
%% point: 单摆
%% method: 无
%% score: 2
%% degree: 650
%% duplicateId: 0
%% body:
一物体在某行星表面受到的万有引力是它在地球表面受到的万有引力的 1/4。在地球上走得很准的摆钟搬到此行星上后，此钟的分针走一整圈所经历的时间实际上是\xzanswer{C}
\fourchoices[answer=C]
{1/4 小时}
{1/2 小时}
{2 小时}
{4 小时}

%% answer: C
%% memo:
\memoanswer{
单摆周期公式为 $T = 2\pi \sqrt{\frac{l}{g}}$，因为 $\frac{g'}{g} = \frac{1}{4}$，所以 $\frac{T'}{T} = 2$。在地球表面，分针走一整圈的时间为 $1\Uh$，所以在该行星表面为 $2\Uh$，故 C 正确。
}

\item
%% number: 3
%% paperName: 1990年全国高考第 3 题 2 分
%% typeId: 1
%% type: 单选题
%% chapter: 光学
%% point: 光电效应
%% method: 无
%% score: 2
%% degree: 650
%% duplicateId: 0
%% body:
用绿光照射一光电管，能产生光电效应。欲使光电子从阴极逸出时的最大初动能增大，应\xzanswer{D}
\fourchoices[answer=D]
{改用红光照射}
{增大绿光的强度}
{增大光电管上的加速电压}
{改用紫光照射}

%% answer: D
%% memo:
\memoanswer{
由光电效应方程有 $h\nu - W_0 = \frac{1}{2}mv_0^2$，最大初动能增大与仅频率有关，因 $\nu_{\text{紫}} > \nu_{\text{绿}}$，所以光电子逸出的最大初动能增大，故 D 正确。
}

\item
%% number: 4
%% paperName: 1990年全国高考第 4 题 2 分
%% typeId: 1
%% type: 单选题
%% chapter: 原子和原子核
%% point: 玻尔模型
%% method: 无
%% score: 2
%% degree: 650
%% duplicateId: 0
%% body:
按照玻尔理论，一个氢原子中的电子从一半径为 $r_{a}$ 的圆轨道自发地直接跃迁到一半径为 $r_{b}$ 的圆轨道上，$r_{a} > r_{b}$，在此过程中\xzanswer{C}
\fourchoices[answer=C]
{原子要发出一系列频率的光子}
{原子要吸收一系列频率的光子}
{原子要发出某一频率的光子}
{原子要吸收某一频率的光子}

%% answer: C
%% memo:
\memoanswer{
原子向低轨跃迁时，会发出光子，设原子跃迁发出光子的频率为 $\nu$，有 $h\nu = E_a - E_b$，所以原子发出光子的频率为 $\frac{E_a - E_b}{h}$，故 C 正确。
}

\item
%% number: 5
%% paperName: 1990年全国高考第 5 题 2 分
%% typeId: 1
%% type: 单选题
%% chapter: 电路
%% point: 电容
%% method: 无
%% score: 2
%% degree: 650
%% duplicateId: 0
%% body:
电容器 $C_{1}$、$C_{2}$ 和可变电阻器 $R_{1}$、$R_{2}$ 以及电源 $\varepsilon$ 连接成如图所示的电路。当 $R_{1}$ 的滑动触头在图示位置时，$C_{1}$、$C_{2}$ 的电量相等。要使 $C_{1}$ 的电量大于 $C_{2}$ 的电量，应\xzanswer{D}
\onepicture[w=5cm]{figs/05.png}
\fourchoices[answer=D]
{增大 $R_{2}$}
{减小 $R_{2}$}
{将 $R_{1}$ 的滑动触头向 A 端移动}
{将 $R_{1}$ 的滑动触头向 B 端移动}

%% answer: D
%% memo:
\memoanswer{
因 $C_{1}$ 和 $C_{2}$ 的电压比例等于 $R_{1}$ 的左端与右端电阻的比例，该比例由滑动变阻器的滑动头位置决定，与 $R_{2}$ 无关，故 AB 错误；要使 $C_{1}$ 的电量大于 $C_{2}$ 的电量，应使 $C_{1}$ 的电压大于 $C_{2}$ 的电压，故滑动触头应向 B 端移动，故 D 正确。
}

\item
%% number: 6
%% paperName: 1990年全国高考第 6 题 2 分
%% typeId: 1
%% type: 单选题
%% chapter: 功和能
%% point: 动能定理
%% method: 无
%% score: 2
%% degree: 650
%% duplicateId: 0
%% body:
一质量为 $2\Ukg$ 的滑块，以 $4\Ums$ 的速度在光滑水平面上向左滑行。从某一时刻起，在滑块上作用一向右的水平力。经过一段时间，滑块的速度方向变为向右，大小为 $4\Ums$。在这段时间里水平力做的功为\xzanswer{A}
\fourchoices[answer=A]
{0}
{$8\UJ$}
{$16\UJ$}
{$32\UJ$}

%% answer: A
%% memo:
\memoanswer{
对滑块的运动过程分析，由动能定理得 $W=\frac{1}{2}mv_{2}^{2}-\frac{1}{2}mv_{1}^{2}=0$，故 A 正确。
}

\item
%% number: 7
%% paperName: 1990年全国高考第 7 题 2 分
%% typeId: 1
%% type: 单选题
%% chapter: 交流电
%% point: LC振荡回路
%% method: 无
%% score: 2
%% degree: 650
%% duplicateId: 0
%% body:
已知 LC 振荡电路中电容器极板 1 上的电量随时间变化的曲线如下图所示。则\xzanswer{D}
\onepicture[w=7cm]{figs/07.png}
\fourchoices[answer=D]
{a、c 两时刻电路中电流最大，方向相同}
{a、c 两时刻电路中电流最大，方向相反}
{b、d 两时刻电路中电流最大，方向相同}
{b、d 两时刻电路中电流最大，方向相反}

%% answer: D
%% memo:
\memoanswer{
由 $i=\frac{\Delta q}{\Delta t}$ 知，$i$ 为 $q-t$ 图象切线的斜率，所以，a、c 两时刻电流为 0，b、d 两时刻电流最大，且方向相反，故 D 正确。
}

\item
%% number: 8
%% paperName: 1990年全国高考第 8 题 2 分
%% typeId: 1
%% type: 单选题
%% chapter: 功和能
%% point: 动能定理
%% method: 无
%% score: 2
%% degree: 550
%% duplicateId: 0
%% body:
三个相同的带电小球 1、2、3，在重力场中从同一高度由静止开始落下，其中小球 1 通过一附加的水平方向匀强电场，小球 2 通过一附加的水平方向匀强磁场。设三个小球落到同一高度时的动能分别为 $E_{1}$、$E_{2}$ 和 $E_{3}$，忽略空气阻力，则\xzanswer{B}
\fourchoices[answer=B]
{$E_{1} = E_{2} = E_{3}$}
{$E_{1} > E_{2} = E_{3}$}
{$E_{1} < E_{2} = E_{3}$}
{$E_{1} > E_{2} > E_{3}$}

%% answer: B
%% memo:
\memoanswer{
带电小球运动时，因电场力做功，增加了小球的动能，而洛伦兹力不做功，所以 $E_{1} > E_{2} = E_{3}$，故 B 正确。
}

\item
%% number: 9
%% paperName: 1990年全国高考第 9 题 2 分
%% typeId: 1
%% type: 单选题
%% chapter: 气体性质
%% point: 阿伏伽德罗常数
%% method: 无
%% score: 2
%% degree: 750
%% duplicateId: 0
%% body:
从下列哪一组数据可以算出阿伏伽德罗常数\xzanswer{D}
\fourchoices[answer=D]
{水的密度和水的摩尔质量}
{水的摩尔质量和水分子的体积}
{水分子的体积和水分子的质量}
{水分子的质量和水的摩尔质量}

%% answer: D
%% memo:
\memoanswer{
A 项可求出水的摩尔体积，故 A 错误；C 项可求出水分子的密度，故 C 错误；D 项可求出阿伏伽德罗常数 $N_{A} = \frac{M}{m}$，故 D 正确，B 错误。
}

\item
%% number: 10
%% paperName: 1990年全国高考第 10 题 2 分
%% typeId: 1
%% type: 单选题
%% chapter: 力物体的平衡
%% point: 力矩平衡
%% method: 极值
%% score: 2
%% degree: 450
%% duplicateId: 0
%% body:
一均匀的直角三角形木板 ABC，可绕垂直纸面通过 C 点的水平轴转动，如图。现用一始终沿直角边 AB 的、作用于 A 点的力 $F$，使 BC 边缓慢地由水平位置转至竖直位置。在此过程中，力 $F$ 的大小随 $\alpha$ 角变化的图线是\xzanswer{D}
\onepicture[w=3.6cm]{figs/10.png}
\fourchoices[answer=D, ispicture=true, h=2.6cm]
{figs/10a.png}
{figs/10b.png}
{figs/10c.png}
{figs/10d.png}

%% answer: D
%% memo:
\memoanswer{
由题意知，对模板受力分析，如图所示，设重心为 O，由力矩平衡得 $F \cdot CB = mg \cdot OC \cos(\alpha + \beta)$（O 为三角形 ABC 的重心），$\beta$ 为 $\angle OCB$，$\beta$ 不变，$\alpha$ 越大，$(\alpha + \beta)$ 越小，$F$ 越小，并且不是线性变化。

\onepicture[w=4cm]{figs/10_an.png}

当重心 O 与 C 的连线在竖直线上时，$F = 0$，然后 $F$ 非线性增大，故 D 正确。
}

\item
%% number: 11
%% paperName: 1990年全国高考第 11 题 2 分
%% typeId: 1
%% type: 单选题
%% chapter: 电磁感应
%% point: 二级电磁感应
%% method: 无
%% score: 2
%% degree: 450
%% duplicateId: 0
%% body:
图中 T 是绕有两组线圈的闭合铁心，线圈的绕向如图所示，D 是理想的二极管，金属棒 ab 可在两条平行的金属导轨上沿导轨滑行，磁场方向垂直纸面向里。若电流计 G 中有电流通过，则 ab 棒的运动可能是\xzanswer{C}
\onepicture[w=5cm]{figs/11.png}
\fourchoices[answer=C]
{向左匀速运动}
{向右匀速运动}
{向左匀加速运动}
{向右匀加速运动}

%% answer: C
%% memo:
\memoanswer{
ab 棒匀速运动时，产生恒定的电流，故左线圈无感应电流，故 AB 错误；ab 棒向左匀加速运动时，ab 中电流 $I_{1}$ 增大，方向如图所示，$I_{1}$ 的磁场方向如图中 $B_{1}$，则左线圈感应电流方向的磁场方向如图中 $B_{2}$，感应电流的方向如图中 $I_{2}$，则二极管为正向接法，G 中有电流，故 C 正确；ab 棒向右匀加速运动时，则左线圈感应电动势方向使二极管为反向接法，所以 G 中无电流，故 D 错误。

\onepicture[w=5.5cm]{figs/11_an.png}
}

\item
%% number: 12
%% paperName: 1990年全国高考第 12 题 2 分
%% typeId: 1
%% type: 单选题
%% chapter: 磁场
%% point: 洛伦兹力
%% method: 无
%% score: 2
%% degree: 650
%% duplicateId: 0
%% body:
一个带电粒子，沿垂直于磁场的方向射入一匀强磁场，粒子的一段径迹如图所示。径迹上的每一小段都可近似看成圆弧。由于带电粒子使沿途的空气电离，粒子的能量逐渐减小（带电量不变）。从图中情况可以确定\xzanswer{B}
\onepicture[w=2.6cm]{figs/12.png}
\fourchoices[answer=B]
{粒子从 a 到 b，带正电}
{粒子从 b 到 a，带正电}
{粒子从 a 到 b，带负电}
{粒子从 b 到 a，带负电}

%% answer: B
%% memo:
\memoanswer{
由题意知，因能量逐渐减小，所以速度逐渐减小，由 $R=\frac{mv}{qB}$ 知半径逐渐减小，所以粒子从 b 到 a，由左手定则知，粒子带正电，故 B 正确。
}

\item
%% number: 13
%% paperName: 1990年全国高考第 13 题 2 分
%% typeId: 1
%% type: 单选题
%% chapter: 力物体的平衡
%% point: 物体系平衡
%% method: 整体与隔离
%% score: 2
%% degree: 650
%% duplicateId: 0
%% body:
如图，在粗糙的水平面上放一三角形木块 a，若物体 b 在 a 的斜面上匀速下滑，则\xzanswer{A}
\onepicture[w=5cm]{figs/13.png}
\fourchoices[answer=A]
{a 保持静止，而且没有相对于水平面运动的趋势}
{a 保持静止，但有相对于水平面向右运动的趋势}
{a 保持静止，但有相对于水平面向左运动的趋势}
{因未给出所需数据，无法对 a 是否运动或有无运动趋势作出判断}

%% answer: A
%% memo:
\memoanswer{
对 b 受力分析：支持力 $N = mg\cos\alpha$，摩擦力 $f = mg\sin\alpha$（因匀速运动，受力平衡）；a 受力的水平分量：压力的 $N_x = mg\cos\alpha\cdot\sin\alpha$，摩擦力的 $f_x = mg\sin\alpha\cdot\cos\alpha$，两力水平分量大小相等，方向相反，所以没有运动趋势，故 A 正确。
}

\item
%% number: 14
%% paperName: 1990年全国高考第 14 题 3 分
%% typeId: 1
%% type: 单选题
%% chapter: 光学
%% point: 光的电磁说
%% method: 无
%% score: 3
%% degree: 650
%% duplicateId: 0
%% body:
设 $\lambda_{1}$、$\lambda_{2}$ 是两种单色可见光 1、2 在真空中的波长。若 $\lambda_{1} > \lambda_{2}$，则这两种单色光相比\xzanswer{AC}
\fourchoices[answer=AC]
{单色光 1 的频率较小}
{玻璃对单色光 1 的折射率较大}
{在玻璃中，单色光 1 的传播速度较大}
{单色光 1 的光子的能量较大}

%% answer: AC
%% memo:
\memoanswer{
由 $c = \lambda \nu$ 得 $\nu = \frac{c}{\lambda}$，因 $\lambda_1 > \lambda_2$，所以 $\nu_1 < \nu_2$，故 A 正确；频率越大，折射率越大，所以可见光 1 的折射率小，故 B 错误；由 $n = \frac{c}{v}$ 得 $v = \frac{c}{n}$，因 $n_1 < n_2$，所以 $v_1 > v_2$，故 C 正确；由 $E = h \nu$ 知，因 $\nu_1 < \nu_2$，所以 $E_1 < E_2$，故 D 错误。
}

\item
%% number: 15
%% paperName: 1990年全国高考第 15 题 3 分
%% typeId: 2
%% type: 多选题
%% chapter: 电场
%% point: 电势能电势差电场力做功
%% method: 无
%% score: 3
%% degree: 650
%% duplicateId: 0
%% body:
一带电粒子射入一固定在 O 点的点电荷的电场中，粒子运动轨迹如图中虚线 abc 所示。图中实线是同心圆弧，表示电场的等势面。不计重力，可以判断\xzanswer{ABD}
\onepicture[w=2.6cm]{figs/15.png}
\fourchoices[answer=ABD]
{此粒子一直受到静电排斥力作用}
{粒子在 b 点的电势能一定大于在 a 点的电势能}
{粒子在 b 点的速度一定大于在 a 点的速度}
{粒子在 a 点和 c 点的速度大小一定相等}

%% answer: ABD
%% memo:
\memoanswer{
由轨迹判断，粒子受斥力作用，故 A 正确；从 a 到 b 电场力做负功，所以电势能增大，故 B 正确；从 a 到 b，电势能增大，则动能减少，故 C 错误；因电势能与动能之和恒定，在同一个等势面上，电势能相等，则动能相等，速度大小相等，故 D 正确。
}

\item
%% number: 16
%% paperName: 1990年全国高考第 16 题 3 分
%% typeId: 2
%% type: 多选题
%% chapter: 曲线运动
%% point: 平抛运动
%% method: 无
%% score: 3
%% degree: 450
%% duplicateId: 0
%% body:
向空中发射一物体，不计空气阻力。当此物体的速度恰好沿水平方向时，物体炸裂成 a、b 两块，若质量较大的 a 块的速度方向仍沿原来的方向，则\xzanswer{CD}
\fourchoices[answer=CD]
{b 的速度方向一定与原速度方向相反}
{从炸裂到落地的这段时间里，a 飞行的水平距离一定比 b 的大}
{a、b 一定同时到达水平地面}
{在炸裂过程中，a、b 受到的爆炸力的冲量大小一定相等}

%% answer: CD
%% memo:
\memoanswer{
物体的速度恰好沿水平方向时，速度的大小为 $v$，由动量守恒得 $mv = m_a v_a + m_b v_b$，如果 $v_a < v$，则 b 的速度方向与原速度方向相同，故 A 错误；不能确定 $v_a$ 和 $v_b$ 哪个大，所以从炸裂到落地的这段时间里，a 飞行的水平距离不一定比 b 的大，故 B 错误；由 $h = \frac{1}{2}gt^2$ 知 h 相等时，t 相等，a、b 一定同时到达水平地面，故 C 正确；由牛顿第三定律知，爆炸力 F 大小相等，同时存在，同时消失，由 $I = F\Delta t$ 知，在炸裂过程中，a、b 受到的爆炸力的冲量大小一定相等，故 D 正确。
}

\item
%% number: 17
%% paperName: 1990年全国高考第 17 题 3 分
%% typeId: 2
%% type: 多选题
%% chapter: 曲线运动
%% point: 人造地球卫星
%% method: 无
%% score: 3
%% degree: 650
%% duplicateId: 0
%% body:
假如一作圆周运动的人造地球卫星的轨道半径增大到原来的 2 倍，仍做圆周运动，则\xzanswer{CD}
\fourchoices[answer=CD]
{根据公式 $v = \omega r$，可知卫星运动的线速度将增大到原来的 2 倍}
{根据公式 $F = \frac{mv^{2}}{r}$，可知卫星所需的向心力将减小到原来的 $\frac{1}{2}$}
{根据公式 $F = \frac{GMm}{r^{2}}$，可知地球提供的向心力将减小到原来的 $\frac{1}{4}$}
{根据（B）和（C）中给出的公式，可知卫星运动的线速度将减小到原来的 $\frac{\sqrt{2}}{2}$}

%% answer: CD
%% memo:
\memoanswer{
由万有引力的公式得 $F = G \frac{Mm}{r^2}$，因公式中除 $r$ 外，其他物理量都不变，所以万有引力减小到原来的 $\frac{1}{4}$，又向心力由万有引力提供，所以向心力减小到原来的 $\frac{1}{4}$，故 C 正确；由 $G \frac{Mm}{R^2} = m \frac{v^2}{r}$ 得 $v = \sqrt{\frac{GM}{r}}$，所以 $\frac{v_2}{v_1} = \sqrt{\frac{r_1}{r_2}} = \sqrt{\frac{1}{2}}$，故 D 正确；A 错误的原因：$\omega$ 是变化的，$v$ 不能由 $r$ 唯一确定，故 B 错误。
}

\item
%% number: 18
%% paperName: 1990年全国高考第 18 题 3 分
%% typeId: 2
%% type: 多选题
%% chapter: 原子和原子核
%% point: 天然放射性
%% method: 无
%% score: 3
%% degree: 650
%% duplicateId: 0
%% body:
\ce{^{232}_{90}Th}（钍）经过一系列 $\alpha$ 和 $\beta$ 衰变，成为 \ce{^{208}_{82}Pb}（铅），\xzanswer{ABD}
\fourchoices[answer=ABD]
{铅核比钍核少 8 个质子}
{铅核比钍核少 16 个中子}
{共经过 4 次 $\alpha$ 衰变和 6 次 $\beta$ 衰变}
{共经过 6 次 $\alpha$ 衰变和 4 次 $\beta$ 衰变}

%% answer: ABD
%% memo:
\memoanswer{
核反应前后满足质量数和核电荷数守恒，质子数的差值为 $90 - 82 = 8$，故 A 正确；中子数的差值为 $(232 - 208) - 8 = 16$，故 B 正确；$\alpha$ 衰变次数 $\frac{232 - 208}{4} = 6$，故 C 错误；$\beta$ 衰变次数 $\frac{90 - 82 - 2 \times 6}{-1} = 4$，故 D 正确。
}

\item
%% number: 19
%% paperName: 1990年全国高考第 19 题 3 分
%% typeId: 1
%% type: 单选题
%% chapter: 气体性质
%% point: 热力学第一定律
%% method: 无
%% score: 3
%% degree: 650
%% duplicateId: 0
%% body:
一定量气体可经不同的过程从状态（$p_{1}$、$V_{1}$、$T_{1}$）变到状态（$p_{2}$、$V_{2}$、$T_{2}$），已知 $T_{2} > T_{1}$，则在这些过程中\xzanswer{D}
\fourchoices[answer=D]
{气体一定都从外界吸收热量}
{气体和外界交换的热量都是相等的}
{外界对气体所做的功都是相等的}
{气体内能的变化量都是相等的}

%% answer: D
%% memo:
\memoanswer{
由热力学第一定律知 $\Delta U = Q + W$，虽然温度升高，$\Delta U$ 为正，但功 W 的正、负不能确定，所以 Q 的正、负不能确定，故 A 错误；同理 BC 错误；因不同变化过程中，温度的变化都相等，所以气体内能的变化量都相等，故 D 正确。
}

\item
%% number: 20
%% paperName: 1990年全国高考第 20 题 3 分
%% typeId: 2
%% type: 多选题
%% chapter: 电磁感应
%% point: 电磁感应中图象
%% method: 无
%% score: 3
%% degree: 450
%% duplicateId: 0
%% body:
一闭合线圈固定在垂直于纸面的匀强磁场中。设向里为磁感应强度 $B$ 的正方向，线圈中的箭头为电流 $i$ 的正方向（如图（a）所示）。已知线圈中感生电流 $i$ 随时间而变化的图象如图（b）所示。则磁感应强度 $B$ 随时间而变化的图象可能是\xzanswer{CD}
\onepicture[w=8cm]{figs/20.png}
\fourchoices[answer=CD, ispicture=true, h=1.8cm]
{figs/20a.png}
{figs/20b.png}
{figs/20c.png}
{figs/20d.png}

%% answer: CD
%% memo:
\memoanswer{
从题图知：$0\sim0.5\Us$，感应电流为负；$0.5\sim1\Us$，感应电流为正；$1\sim1.5\Us$，感应电流为正；$1.5\sim2\Us$，感应电流为负。A 项：$0\sim0.5\Us$，B 垂直纸面向外，增加，感应电流产生的磁场垂直纸面向里，感应电流同图 1 为正，故 A 错误；B 项：$0\sim0.5\Us$，B 垂直纸面向外，减小，感应电流为负；$0.5\sim1\Us$，B 垂直纸面向里，增大，感应电流为负，故 B 错误；C 项：$0\sim0.5\Us$，B 垂直纸面向里，增大，感应电流为负；$0.5\sim1\Us$，B 垂直纸面向里，减小，感应电流为正，故 C 正确；D 项：感应电流方向同 C，故 D 正确。
}

\item
%% number: 21
%% paperName: 1990年全国高考第 21 题 3 分
%% typeId: 1
%% type: 单选题
%% chapter: 电场
%% point: 带电粒子的平衡
%% method: 整体与隔离
%% score: 3
%% degree: 450
%% duplicateId: 0
%% body:
用轻质细线把两个质量未知的小球悬挂起来，如右图所示。今对小球 a 持续施加一个向左偏下 $30^{\circ}$ 的恒力，并对小球 b 持续施加一个向右偏上 $30^{\circ}$ 的同样大的恒力，最后达到平衡。表示平衡状态的图可能是\xzanswer{A}
\onepicture[w=2.2cm]{figs/21.png}
\fourchoices[answer=A, ispicture=true, h=3cm]
{figs/21a.png}
{figs/21b.png}
{figs/21c.png}
{figs/21d.png}

%% answer: A
%% memo:
\memoanswer{
对 a、b 整体分析，a、b 受力平衡，两个 F 相互抵消，只受重力和绳的拉力作用，二力平衡，在同一直线上，悬线竖直，故 A 正确。
}

\item
%% number: 22
%% paperName: 1990年全国高考第 22 题 3 分
%% typeId: 3
%% type: 填空题
%% chapter: 电磁感应
%% point: 自感与互感，涡流
%% method: 无
%% score: 3
%% degree: 650
%% duplicateId: 0
%% body:
图为一演示实验电路图，图中 L 是一带铁心的线圈，A 是一灯泡，电键 K 处于闭合状态，电路是接通的。现将电键 K 打开，则在电路切断的瞬间，通过灯泡 A 的电流方向是从 \tkanswer{a} 端到 \tkanswer{b} 端。这个实验是用来演示 \tkanswer{自感} 现象的。
\onepicture[align=right, w=4cm]{figs/22.png}

%% answer: a，b，自感
\jdanswer{
a，b，自感
}

%% memo:
\memoanswer{
电感线圈 L 上电流从 b 到 a，断电自感，L 上的感应电流方向与原来相同，该电流与灯 A 组成回路，所以灯 A 上电流从 a 到 b。
}

\item
%% number: 23
%% paperName: 1990年全国高考第 23 题 3 分
%% typeId: 3
%% type: 填空题
%% chapter: 曲线运动
%% point: 竖直平面圆周运动
%% method: 无
%% score: 3
%% degree: 650
%% duplicateId: 0
%% body:
一轻绳上端固定，下端连一质量为 $0.05\Ukg$ 的小球。若小球摆动过程中轻绳偏离竖直线的最大角度为 $60^{\circ}$，则小球经过最低点时绳中张力等于 \tkanswer{$1\UN$}。（$g$ 取 $10\Umsq$）

%% answer: 1 N
\jdanswer{
$1\UN$
}

%% memo:
\memoanswer{
小球运动的过程中，只有重力做功，由机械能守恒得 $mgL(1-\cos60^\circ)=\frac{1}{2}mv^2$，在最低点，由牛顿第二定律得 $T-mg=m\frac{v^2}{L}$，解得 $T=2mg=1.0\UN$。
}

\item
%% number: 24
%% paperName: 1990年全国高考第 24 题 3 分
%% typeId: 4
%% type: 实验题
%% chapter: 电路
%% point: 测电动势与内阻
%% method: 无
%% score: 3
%% degree: 650
%% duplicateId: 0
%% body:
用伏安法测电阻的实验中，按实验要求选用的电压表的最小分度为 0.1 伏，电流表的最小分度为 0.02 安。某学生记录的各组数据如下表所示：

\begin{center}
\begin{tabular}{|c|c|c|c|c|c|}
\hline
组次 & 1 & 2 & 3 & 4 & 5 \\
\hline
$U$（V） & 0.81 & 1.21 & 1.7 & 1.79 & 2.51 \\
\hline
$I$（A） & 0.16 & 0.242 & 0.338 & 0.422 & 0.504 \\
\hline
\end{tabular}
\end{center}

在这五组数据中，有效数字位数不符合要求的是第 \tkanswer{1、3} 组，数据有差错的是第 \tkanswer{4} 组。

%% answer: 1、3；4
\jdanswer{
1、3；4
}

%% memo:
\memoanswer{
电压表读到 $0.01\UV$，第 3 组位数不够；电流表读到 $0.001\UA$，第 1 组位数不够。由 $R = \frac{U}{I}$ 知，第 2 组为 $R = 5.0\UO$，第 4 组为 $R = 4.2\UO$，第 5 组为 $R = 5.0\UO$，故第 4 组数有差错。
}

\item
%% number: 25
%% paperName: 1990年全国高考第 25 题 3 分
%% typeId: 3
%% type: 填空题
%% chapter: 交流电
%% point: 交流电
%% method: 无
%% score: 3
%% degree: 650
%% duplicateId: 0
%% body:
额定电压均为 $220\UV$ 的三个相同的灯泡，按星形接法连接在线电压为 $380\UV$、相电压为 $220\UV$ 的三相电路上。如果电路的中性线断了，又将一个灯泡从电路中取下，这时电路中每个灯泡两端的电压为 \tkanswer{$190\UV$}。

%% answer: 190 V
\jdanswer{
$190\UV$
}

%% memo:
\memoanswer{
如果电路的中性线断了，又将一个灯泡从电路中取下，两个灯泡串联接在 $380\UV$ 电压上，所以 $U_{1} + U_{2} = 380\UV$，故 $U_{1} = U_{2} = 190\UV$。
}

\item
%% number: 26
%% paperName: 1990年全国高考第 26 题 3 分
%% typeId: 3
%% type: 填空题
%% chapter: 机械波
%% point: 波的图象
%% method: 无
%% score: 3
%% degree: 750
%% duplicateId: 0
%% body:
如图所示是一列简谐波在 $t=0$ 时的波动图象。波的传播速度为 $2\Ums$，则从 $t=0$ 到 $t=2.5\Us$ 的时间内，质点 M 通过的路程是 \tkanswer{$2.5\Um$}，位移是 \tkanswer{$0\Um$}。
\onepicture[align=right, w=6cm]{figs/26.png}

%% answer: 2.5 m；0
\jdanswer{
$2.5\Um$；$0$
}

%% memo:
\memoanswer{
由图像知 $\lambda = 0.4\Um$，$A = 5\Ucm$，由 $v = \frac{\lambda}{T}$ 得 $T = \frac{\lambda}{v} = \frac{0.4}{2}\Us = 0.2\Us$，又 $t = 2.5\Us = 12\frac{1}{2}T$，故质点 M 通过的路程为 $s = 12.5 \times 4A = 2.5\Um$，位移为 $x = 0$。
}

\item
%% number: 27
%% paperName: 1990年全国高考第 27 题 3 分
%% typeId: 4
%% type: 实验题
%% chapter: 电路
%% point: 多用表的使用
%% method: 无
%% score: 3
%% degree: 650
%% duplicateId: 0
%% body:
用万用表欧姆挡（$\times100$）测试三只晶体二极管，其结果依次如图①、②、③所示。由图可知，图 \tkanswer{②} 中的二极管是好的，该二极管的正极是 \tkanswer{a} 端。
\onepicture[align=right, w=6cm]{figs/27.png}

%% answer: ②；a
\jdanswer{
②；a
}

%% memo:
\memoanswer{
图②中的二极管是好的，因二极管具有单向导电性，该二极管的正极是 a 端，因欧姆表黑表笔（表盘“$-$”）接的是内电池的正极，二极管接 a 端时是导通的（电阻小接近 0）。
}

\item
%% number: 28
%% paperName: 1990年全国高考第 28 题 3 分
%% typeId: 4
%% type: 实验题
%% chapter: 直线运动
%% point: 游标卡尺和螺旋测微仪
%% method: 无
%% score: 3
%% degree: 650
%% duplicateId: 0
%% body:
图中给出的是用螺旋测微器测量一小钢球的直径时的示数，此读数应是 \tkanswer{$8.600\Umm$}。
\onepicture[align=right, w=4cm]{figs/28.png}

%% answer: 8.600 mm
\jdanswer{
$8.600\Umm$
}

%% memo:
\memoanswer{
固定刻度为 $8.5\Umm$，可动刻度为 $10.0 \times 0.01 = 0.100\Umm$，故读数为 $8.5\Umm + 0.100\Umm = 8.600\Umm$。
}

\item
%% number: 29
%% paperName: 1990年全国高考第 29 题 3 分
%% typeId: 4
%% type: 实验题
%% chapter: 电路
%% point: 电表改装
%% method: 无
%% score: 3
%% degree: 450
%% duplicateId: 0
%% body:
如图所示是把量程为 $3\UmA$ 的电流表改装成欧姆表的结构示意图，其中电池电动势 $\varepsilon = 1.5\UV$。经改装后，若将原电流表 $3\UmA$ 刻度处的刻度值定为零位置，则 $2\UmA$ 刻度处应标 \tkanswer{$250\UO$}，$1\UmA$ 刻度处应标 \tkanswer{$1000\UO$}。
\onepicture[align=right, w=4cm]{figs/29.png}

%% answer: 250 Ω；1000 Ω
\jdanswer{
$250\UO$；$1000\UO$
}

%% memo:
\memoanswer{
由闭合电路的欧姆定律知 $I = 3\UmA = \frac{E}{R_{\text{内}}}$，$R_{\text{内}} = \frac{1.5}{3}\UkO = 0.5\UkO = 500\UO$，又 $I_1 = 2\UmA = \frac{E}{R_{\text{内}} + R_{x_1}}$，解得 $R_{x1} = \frac{1.5}{2 \times 10^{-3}} - 500 = 250\UO$；又 $I_3 = 1\UmA = \frac{E}{R_{\text{内}} + R_{x_2}}$，解得 $R_{x2} = \frac{1.5}{1 \times 10^{-3}} - 500 = 1000\UO$。
}

\item
%% number: 30
%% paperName: 1990年全国高考第 30 题 5 分
%% typeId: 7
%% type: 作图题
%% chapter: 光学
%% point: 透镜
%% method: 无
%% score: 5
%% degree: 650
%% duplicateId: 0
%% body:
图中 MN 是薄透镜的主轴，S 是发光点，S$'$ 是它的像点。

\begin{enumerate}
	\item
	用作图法求出薄透镜的位置，标在图上。
	\item
	分别作光路图求出两个焦点的位置，标在图上。再标明透镜的类别。
\end{enumerate}
\drawpicanswer{\onepicture[align=right, w=5cm]{figs/30.png}}{\onepicture[align=right, w=5cm]{figs/30_an.png}}

%% answer: （1）如图；（2）如图
\jdanswer{
\begin{enumerate}
	\item
	如图

	\item
	如图
\end{enumerate}
}

%% memo:
\memoanswer{
参考解答如图。

评分标准：本题共 5 分。用作图法找到光心 O 给 1 分，找到一个焦点再给 2 分，找到第二个焦点再给 1 分，光路图完整而正确的再给 1 分。（光心 O 和两个焦点都必须在 MN 轴上，它们的横坐标分别是 4.8 到 5.2 之间、7.1 到 7.9 之间和 2.1 到 2.9 之间，否则按评分标准扣除该项的得分。凡不用光路图得到的结果都不给分。光路图不完整或透镜类别、光线箭头、虚实线等任一部分有错的，都不给最后那 1 分。）
}

\item
%% number: 31
%% paperName: 1990年全国高考第 31 题 6 分
%% typeId: 6
%% type: 计算题
%% chapter: 气体性质
%% point: 气体连接体
%% method: 无
%% score: 6
%% degree: 650
%% duplicateId: 0
%% body:
用销钉固定的活塞把水平放置的容器分隔成 A、B 两部分，其体积之比 $V_{A}:V_{B} = 2:1$，如图所示。起初 A 中有温度为 $127\UdC$、压强为 $1.8\times10^{5}\UPa$ 的空气，B 中有温度 $27\UdC$、压强为 $1.2\times10^{5}\UPa$ 的空气。拔出销钉，使活塞可以无摩擦地移动（不漏气）。由于容器壁缓慢导热，最后气体都变到室温 $27\UdC$，活塞也停住，求最后 A 中气体的压强。
\onepicture[align=right, w=4.5cm]{figs/31.png}

%% answer: $1.3\times10^{5}\UPa$
\jdanswer{
$1.3\times10^{5}\UPa$
}

%% memo:
\memoanswer{
开始，A 和 B 中气体的压强、体积、温度分别为 $p_{A}$、$V_{A}$、$T_{A}$ 和 $p_{B}$、$V_{B}$、$T_{B}$，且 $V_{A}=2V_{B}$。最后两部分气体的压强都是 $p$，温度都是 $T$，体积分别是 $V_{A}'$、$V_{B}'$。由气态方程，有

$\frac{p_{A}V_{A}}{T_{A}}=\frac{pV_{A}'}{T}$（1）

$\frac{p_{B}V_{B}}{T_{B}}=\frac{pV_{B}'}{T}$（2）

又 $V_{A}' + V_{B}' = V_{A} + V_{B}$（3）

代入已知条件，解得 $p = T\left(\frac{2p_{A}}{3T_{A}} + \frac{p_{B}}{3T_{B}}\right) = 1.3 \times 10^{5}\UPa$。

评分标准：全题 6 分。正确列出（1）、（2）两式给 3 分（仅仅写出一般的气态方程，而没有体现两边气体末态的压强相等、温度相等的，不给这 3 分）。正确列出（3）式再给 1 分。正确解出压强 p 再给 2 分（数值、单位各占 1 分）。
}

\item
%% number: 32
%% paperName: 1990年全国高考第 32 题 7 分
%% typeId: 6
%% type: 计算题
%% chapter: 电磁感应
%% point: 电磁感应电路分析
%% method: 无
%% score: 7
%% degree: 650
%% duplicateId: 0
%% body:
固定在匀强磁场中的正方形导线框 abcd，各边长为 $l$，其中 ab 是一段电阻为 $R$ 的均匀电阻丝，其余三边均为电阻可忽略的铜线。磁场的磁感应强度为 $B$，方向垂直纸面向里。现有一与 ab 段材料、粗细、长度都相同的电阻丝 PQ 架在导线框上（如图），以恒定的速度 $v$ 从 ad 滑向 bc。当 PQ 滑过 $\frac{l}{3}$ 的距离时，通过 aP 段电阻丝的电流强度是多大？方向如何？
\onepicture[align=right, w=3.6cm]{figs/32.png}

%% answer: $I_{aP} = \frac{6Blv}{11R}$，由 P 流向 a
\jdanswer{
$I_{aP} = \frac{6Blv}{11R}$，由 P 流向 a
}

%% memo:
\memoanswer{
把 PQ 作为电源，内阻为 $R$，电动势为 $E$。

$E = Blv$（1）

外电路是一并联电路，$r_{\text{外}}=\frac{\frac{1}{3}R\times\frac{2}{3}R}{\frac{1}{3}R+\frac{2}{3}R}=\frac{2}{9}R$（2）

电路总电阻 $r_{\text{总}}=R+r_{\text{外}}=\frac{11}{9}R$（3）

总电流 $I=\frac{E}{r_{\text{总}}}=\frac{9Blv}{11R}$（4）

由分流得 $I_{aP}=\frac{2}{3}I$（5）

aP 段中电流的大小 $I_{aP} = \frac{6Blv}{11R}$，由 P 流向 a。

评分标准：全题 7 分。正确列出（1）式得 1 分。正确得出（2）、（3）、（4）、（5）式各得 1 分。正确得出 aP 段中电流的大小和流向再各得 1 分。
}

\item
%% number: 33
%% paperName: 1990年全国高考第 33 题 8 分
%% typeId: 6
%% type: 计算题
%% chapter: 力物体的平衡
%% point: 力矩平衡
%% method: 无
%% score: 8
%% degree: 450
%% duplicateId: 0
%% body:
质量 $m = 2.0\Ukg$ 的小铁块静止于水平导轨 AB 的 A 端。导轨及支架 ABCD 形状及尺寸如上右图，它只能绕通过支架 D 点的垂直于纸面的水平轴转动，其重心在图中的 O 点，质量 $M = 4.0\Ukg$。现用一细线沿导轨拉铁块，拉力 $F = 12\UN$。铁块和导轨之间的摩擦系数 $\mu = 0.50$。重力加速度 $g = 10\Umsq$。从铁块运动时起，导轨（及支架）能保持静止的最长时间是多少？
\onepicture[align=right, w=4.5cm]{figs/33.png}

%% answer: $1.0\Us$
\jdanswer{
$1.0\Us$
}

%% memo:
\memoanswer{
导轨刚要不能维持平衡时，C 端受的力为零，此时导轨（及支架）受四个力：滑块对导轨的压力 $N = mg$，竖直向下；滑块对导轨的摩擦力 $f = \mu mg = 10\UN$，方向向右；重力 $Mg$，作用在 O 点，方向竖直向下；轴作用于 D 端的力。

设此时铁块走过路程 $s$，根据有轴物体平衡条件及图中尺寸，有

$0.1Mg + mg(0.7 - s) = f \times 0.8 = 0.8\mu mg$

$40 \times 0.1 + 20(0.7 - s) = 10 \times 0.8$（1）

解得 $s = 0.50\Um$。

铁块受的摩擦力 $f = 10\UN$，向左，由牛顿第二定律得

$F - f = ma, \quad 12 - 10 = 2a$（2）

解得 $a = 1.0\Umsq$。

将 s 和 a 代入 $t = \sqrt{\frac{2s}{a}} = 1.0\Us$。

评分标准：全题 8 分。正确列出（1）式得 4 分，解得 s = 0.50 米再得 1 分，共计 5 分。凡因力的分析、力矩的大小和转向等而导致（1）式错误的就不给这 5 分，但（1）式正确而 s 算错的给 4 分。正确列出（2）式得 2 分。求出正确结果 t = 1.0 秒再得 1 分。
}

\end{enumerate}

\end{document}
